\section{Arithmetic realization of the local level and weight comparisons}
\label{sec:arithmetic-realization}

The datum $D$ remains fixed throughout this section. Put $E=D.\mathrm{host}.E$,
$F=E^+$, $n=D.R.nV=r+t$, and $K_{i,k}=\operatorname{Frac}(O_{i,k})$.
The index $j$ always lies in $\operatorname{Fin}(D.R.m)$; its cell is at
position $j+1$. The group of finite adelic unitary points is $G_f$, and its
rational subgroup is $\Gamma$. We use multiplicative valuations, so
membership in the maximal ideal means valuation strictly less than $1$.
At a split place $s$, write $v=s.v$, $P=s.w$, and
$\operatorname{loc}_s:G_f\to\mathrm{GL}_n(F_v)$ for the selected component.
The selection matters: the conjugate component is its inverse transpose.

Fix an actual $\texttt{NExpTied}$ object $X$, an admissible twist
$(i,b,k,Q,\psi)$, and its actual cell $M=X.\mathrm{mtV}(i,j+1,b,k,Q,\psi)$.
The cell comes with its coefficient module $W$, representation $\rho$,
level $J$, character $\chi$, realization $\mathrm{real}$, and selected
prime sets $Q_P,D_P$. These are retained data, not constructions from a
dimension formula. For a supplied stage $(B,f,T,w,P_0,x,l)$, write
\[
 \eta=\mu_{0,j}\,\xi(\mu_{\mathrm{Tw}}(v_0))
      =\mu_{0,j}\,\delta_\ell\,\psi_{\mathrm H},
 \qquad
 \delta_\ell=\begin{cases}\delta'&\ell=p_\delta,\\
                          \delta&\ell\ne p_\delta.\end{cases}       \tag{AR.1}
\]
Here $\ell=\operatorname{resChar}(v_0)=D.\mathrm{ntw}.l(i)$, and the
second equality is the actual tuple identity, not an identification of
unrelated characters. The complex characters supplied by HF satisfy the
needed character laws when those HF hypotheses hold. They do not, by
themselves, supply normalized Galois-character compatibility or the
nonempty twist coverage used below.

\paragraph{AR1: exact ramification gives the required bad-place bound.}
Assume the actual $\texttt{IsMu0}$, $\texttt{IsDelta}$,
$\texttt{IsPsi}$ and $\texttt{IsTwist}$ laws in (AR.1).
Outside the two type primes $\mu_{0,j}$ is unramified. The twist character
$\psi_{\mathrm H}$ is ramified exactly at $Q$ and avoids the two auxiliary
rational primes. The selected auxiliary character is unramified outside
the primes over the opposite auxiliary prime, and its inertial image has
order $15$ at every such prime. Then
\[
 \mathrm{badV}(M)\subset g\mathrm{Pl}(D)\cup q\mathrm{Pl}(D,\eta).
                                                               \tag{AR.2}
\]
Moreover the extra finite-order factor $\delta_\ell\psi_{\mathrm H}$ is
unramified at each type prime.

If a prime ideal contains distinct rational primes $p,q$, a Bezout identity
$ap+bq=1$ puts $1$ in that ideal, which is impossible. In the source's
$\texttt{IsTwist}$ condition a ramified prime has prime absolute norm
outside the forbidden rational set. If it also lies above a forbidden
$p$, the norm and $p$ both belong to the prime ideal, so their primality
forces equality and contradicts that exclusion. This proves the claimed
unramifiedness at the auxiliary places, with the literal norm condition.

Now take a bad place and a prime of $E$ above it witnessing badness.
A type prime is already in $g\mathrm{Pl}$. At a prime of $Q$, both
$\mu_{0,j}$ and $\delta_\ell$ are unramified. If their product with
$\psi_{\mathrm H}$ were unramified, cancellation on local units would
make $\psi_{\mathrm H}$ unramified, contrary to its exact ramification.
At an auxiliary prime, $\mu_{0,j}$ and $\psi_{\mathrm H}$ are
unramified; cancellation would make $\delta_\ell$ unramified. Its
inertial image would then have order $1$, contradicting order $15$.
At a type prime, $Q$ is excluded by admissibility and the auxiliary
primes have different rational residue characteristics. Hence both
extra factors are unramified there. These arguments prove (AR.2)
without estimating, enlarging or replacing $Q$.

\paragraph{AR2: selected halves and split-place separation are exact.}
Suppose $P_q$ selects a prime over each prime of $F$. On split primes
assume that the series choice set is exactly $\operatorname{range}P_q$.
For a finite conjugation-stable set $S$ of split primes,
the series half $\operatorname{half}(\operatorname{range}P_q,S)$ equals
the cell half $\operatorname{qpOf}(P_q,S)$.

For a split prime $P$ above $v$, the only primes above $v$ are $P$ and
$\bar P$. Membership in the series half means $P\in S$ and
$P=P_q(v)$. The cell condition is $P\in S$ and either this equality
or exclusion of the other prime from $S$. The second alternative,
when $P_q(v)=\bar P$, is impossible because $S$ is conjugation-stable.
Conversely $P=P_q(v)$ immediately satisfies the cell condition.
This also explains why an arbitrary half cannot be substituted.

The actual sets $Q_P,D_P$ are disjoint, contain no conjugate pair, and
avoid type primes. If selected split primes in $Q_P\cup D_P$ have the
same underlying $v$, they are equal or conjugate. A conjugate pair
within one of the sets is excluded by its defining condition. A
conjugate pair between the sets would put a prime of $Q$ over the
forbidden auxiliary rational prime. A coincidence with a type place
would similarly put the corresponding $g_a$ in a conductor or
auxiliary prime. Thus, at a selected conductor place, it is the only
member of $Q_P\cup D_P$ above that place and no type prime lies above it.
Finally, if $v\in q\mathrm{Pl}(D,\eta)$, choose a prime above it where
$\eta$ is ramified. AR1 shows that it lies in $Q$ or the auxiliary set.
The covering fields of $Q_P,D_P$ select it or its conjugate, producing
the required member of $Q_P\cup D_P$ over precisely $v$.

\paragraph{AR3: tame residue values agree with the specified Artin map.}
Retain the source's published global Artin reciprocity, with its
inertia and inverse-arithmetic-Frobenius convention. Let $\varphi:G_E\to A$
be a homomorphism into a commutative group with open kernel. At a prime
$P$ of residue characteristic $p$, suppose every inertial value has
order dividing a power of an integer $m$ prime to $p$.
Its value on a tame representative depends only on the residue-unit
and the stated exponent. In particular,
\[
 u\equiv1\pmod P\ \Longrightarrow\
 \operatorname{tameVal}(\varphi,P,a,u)=1.                 \tag{AR.3}
\]
If a complex embedding reads $\varphi$ as a Hecke character $\chi_{\mathrm H}$
under the specified Artin map, then for $u\in\mathcal O_E\setminus P$,
\[
 e\!\left(\operatorname{tameVal}(\varphi,P,1,u)\right)
       =\chi_{\mathrm H,P}(u).                          \tag{AR.4}
\]
The notation in (AR.4) includes the literal inclusion of coefficient
units into the fraction field.

The open kernel contains the fixing subgroup of a finite Galois
extension. In a finite local Galois extension the first ramification
quotient acts on a uniformizer through a residue-field unit; each
successive quotient embeds in the additive residue field by the
leading coefficient of $\sigma(\pi)-\pi$. The latter quotients are
$p$-groups and the filtration terminates. Consequently the wild inertia
image in a finite quotient is a $p$-group. Its image under $\varphi$ is
also of $m$-power order, so it is trivial. Tame inertia is read by the
ratios $\sigma(\pi_N)/\pi_N$, for prime-to-$p$ roots of a uniformizer.
Frobenius conjugates its tame generator to its $q$th power, where
$q=\#k_P$. Since the target is commutative, its image has exponent
dividing $q-1$. Thus the ratio for a $(q-1)$st root already determines
the value. Changing the valuation above $P$ conjugates inertia and
does not change a commutative character. This proves independence of
the choices used in $\operatorname{tameVal}$ and proves (AR.3).
These are the usual finite ramification and Hensel facts for the
completed number field; no assertion that wild inertia is trivial
in the absolute Galois group is being made.

The sign in (AR.4) can be checked without a choice of a tame generator.
Choose $\varpi\in P\setminus P^2$, $\pi^{q-1}=\varpi$, and
$t^{q-1}=u$. Arithmetic Frobenius lifts fixing respectively $\pi$ and
$\pi t$ give an inertia element $\tau'^{-1}\tau$. Its ratio on $\pi$
reduces to $u$. Their images under the stated Artin map differ by the
local unit $u$, since the two uniformizer norms are
$(-1)^q\varpi$ and $(-1)^qu\varpi$. For an arbitrary local unit $x$
with the same residue, Hensel lifting gives
$x/u=z^{q-1}$ with $z\equiv1$ modulo the maximal ideal. Lift $z$ to
inertia using the inertia-image clause and multiply by that lift to
the power $q-1$. It changes the Artin image to that of $x$ and leaves
the tame ratio unchanged. Evaluation of the reading identity proves
(AR.4). Replacing the inertia element by its $a$th power gives the
stated exponent-$a$ version.

\paragraph{AR4: the auxiliary side is read with its actual inverse.}
Let $\epsilon_b=1$ for $b=\mathrm{true}$ and $-1$ otherwise.
Assume the literal $\texttt{DeltaRead}(D,cf)$ and
$\texttt{SideRead}(i,b,k,e)$. Then the complex character obtained by
reading $\delta_i^{\epsilon_b}$ is the Hecke character
$cf.\mathrm{delAt}(\ell)$. Hence at $P\in D.\mathrm{ntw}.pD(i)$,
\[
 e\!\left(\operatorname{tameVal}(\delta_i,P,\epsilon_b,u)\right)
       =cf.\mathrm{delAt}(\ell)_P(u).                   \tag{AR.5}
\]

The Hecke character has finite order. RC3 descends it through the
specified Artin quotient to a continuous complex Galois character.
The read character also has open kernel, so is continuous even though
the field embedding $e$ itself is not assumed continuous.
Outside finitely many primes, $\texttt{DeltaRead}$ supplies an element
$a\in L_0$ with
$\iota(a)=cf.\mathrm{delAt}(\ell)(P)^{-1}$ and the prescribed
coefficient equality for $\delta_i(\mathrm{Frob}_P)$.
On the true side, $\texttt{SideRead}$ reads this equality directly.
On the false side the coefficient involution sends $a^{-1}$ to $a$:
the complex values are unitary, so conjugation equals inversion.
The additional inverse in $\delta_i^{\epsilon_b}$ cancels this change.
Thus both sides give the same inverse Frobenius value. RC4 proves
equality of the characters.

The inertia image of $\delta_i$ has order $15$. The auxiliary norm is
$1$ modulo $15$, so that norm is coprime to $15$. Since the norm
is a power of the residue characteristic, that characteristic
does not divide $15$. Apply AR3 with $m=15$, preserving $\epsilon_b$.
On $Q_P$ the analogous application uses the $\ell$-power character
$\psi$ and $\ell\notin P$. AR1 cancels the unramified factors in
(AR.1), so these two formulas read exactly the local component of
$\eta$, not just the separate twist factor.

\paragraph{AR5: the conductor-place subgroup and character are parahoric.}
Let $P\in Q_P\cup D_P$, $s=M.\mathrm{sp}(P)$, and
$\kappa\in G_f$ be trivial outside $s.v$. Then
\[
 \kappa\in M.J
 \quad\Longleftrightarrow\quad
 \operatorname{loc}_s(\kappa)\in\operatorname{ParaRT}(r).
                                                               \tag{AR.6}
\]
Here $\operatorname{ParaRT}(r)$ consists of integral invertible matrices
whose lower-left $t$-by-$r$ block lies in the maximal ideal.
There is a local component $\xi:F_{s.v}^{\times}\to\mathbb C^\times$
of $\eta$ such that, for every allowed reading $e$,
\[
 e(\chi(\kappa))=\xi(\det D_\kappa),                     \tag{AR.7}
\]
where $D_\kappa$ is the lower-right block. This is the literal
$\texttt{LevelPara}$ assertion.

Unfold $M.J$: it imposes global maximal-compact membership, the type
unipotent-block conditions, and parahoric conditions at every prime
of $Q_P,D_P$. AR2 and the support condition make every other local
matrix the identity. The sole remaining condition is the one in
(AR.6). Global maximal-compact membership is equivalent to local
integrality here: the unselected conjugate component is the inverse
transpose, so both a matrix and its inverse are integral exactly when
the selected component is in $\mathrm{GL}_n(\mathcal O_v)$.

Because $\eta$ is tame, compose its $E_P$ component with the split
completion isomorphism $F_v\simeq E_P$ to obtain $\xi$. It is trivial
on $1+\mathfrak m_v$ and gives the required values on $E^\times$.
The block-triangular reduction of an integral invertible matrix has
invertible diagonal blocks, so $\det D_\kappa$ is a unit. Choose its
integral residue representative in $\mathcal O_E$.
The literal product formula for $\chi$ has type factors and one
factor at each member of $Q_P,D_P$. The type factors are $1$ at the
identity. Every other conductor factor has residue $1$, hence equals
$1$ by (AR.3). The unique remaining factor equals
$\xi(\det D_\kappa)$ by (AR.4) or (AR.5).
This proves both the full subgroup equality and its character law.
Together with AR2 it proves $\texttt{paraShape\_export}$.

\paragraph{AR6: a prescribed nontrivial residue character extends to the completion.}
Let $\mathcal O_v\subset F_v$ be the valuation ring and
$\mathrm{red}:\mathcal O_v\to\mathbb F_p$ the specified residue map.
For every nontrivial additive character
$a:\mathbb F_p\to\mathbb C^\times$, there is an additive character
$a_v:F_v\to\mathbb C^\times$ such that
\[
 a_v|_{\mathcal O_v}=a\circ\mathrm{red},\quad
 a_v|_{\mathfrak m_v}=1,\quad
 a_v|_{\mathcal O_v}\ne1.                               \tag{AR.8}
\]

A homomorphism from a subgroup $H$ of an abelian group to
$\mathbb C^\times$ extends to the whole group. Order partial
extensions by inclusion and take a maximal one. If $x$ lies outside
its domain, either no positive multiple of $x$ lies in the domain,
or the least such multiple is $N x$. In the first case give $x$ any
nonzero value. In the second choose an $N$th root of the known value
at $Nx$. The formula on $H+\mathbb Zx$ is well-defined and extends
the homomorphism, contradicting maximality. Apply this to the
additive groups $\mathcal O_v\subset F_v$.
The maximal ideal maps to zero in the residue field. Conversely a
unit cannot map to zero, since ring homomorphisms preserve units;
thus the kernel is exactly the maximal ideal. A residue with
nontrivial $a$ has an integral lift, giving the last assertion.
The character is locally constant because its kernel contains the
open subgroup $\mathfrak m_v$.

\paragraph{AR7: the type character is the actual local component of the base character.}
Retain the exact $\texttt{SideB}$ normalized-character identities
$\texttt{chiD}$, $\texttt{chi}$, $\texttt{eta\_eq}$,
$\texttt{ringSpec\_normalisedCharacter\_eq}$ and
$\texttt{del\_eq}$, as well as $\texttt{ValLL0}$ and
$\texttt{TauTrue}$. Write $\tau_{j+1,a}$ for the single $L_0$-valued
type character supplied by $\texttt{TauTrue}$. Then, for a unit
$z\in\mathbb F_{g_a}^{\times}$ and an integral representative $u$ at
the type prime $P_a$,
\[
 \iota\bigl(\tau_{j+1,a}(z)\bigr)=\mu_{0,j,P_a}(u).
                                                               \tag{AR.9}
\]
This is an equality for this $\tau$, including its orientation.

We retain the algebraic-Hecke-character-to-adic-character existence
input used by $\texttt{GalChar.exists\_isGal}$; finite-order RC3 alone
does not apply to the normalized algebraic block character.
Under that input, construct an adic character for
$\mu_{0,j}\delta_\ell\eta_T^{-1}$ with the source's norm shift.
The normalized compatibility field gives its Frobenius values,
outside a finite set, as exactly those of the actual $\chi_1$.
Use the same reading of $L_0$ on both sides and RC4 to conclude
equality of these characters. This proves the actual
$\texttt{firstTwistCharacter\_eq\_ladic\_heckeCharacter}$ comparison.
No continuity of an isomorphism $\overline{\mathbb Q}_\ell\simeq\mathbb C$
is used; continuity is on the adic side before applying the algebraic
reading.

To apply it at $P_a$, choose a rational prime $p>g_a$ and a coefficient
place above $p$. The \emph{all-place} field $\texttt{Bs.cover}$ supplies
a twist index with $\ell=p\ne g_a$. This step requires a nonempty
twist family. The prime $P_a$ is unramified and of degree one over
$g_a$, so $g_a\in P_a\setminus P_a^2$ is a uniformizer. Apply the
Artin-unit construction of AR3 to a $(g_a-1)$st root of $g_a$.
The actual $\texttt{tau\_galois\_first\_character}$ law identifies the
result with $\tau_{j+1,a}(z)$. In the normalized block character the
auxiliary factor, $\eta_T$, and the norm factor are all trivial on
these local units. What remains is precisely $\mu_{0,j,P_a}(u)$.
This proves (AR.9). Any unramified factor may be multiplied into
$\mu_{0,j}$ without changing (AR.9).

\paragraph{AR8: the type-place subgroup and its nondegenerate character are explicit.}
Fix a type place $s$ above $P_a$. There are characters $a_s,\xi_s$,
chosen \emph{before} the complex reading $e$, such that $a_s$ satisfies
(AR.8), $\xi_s$ is a local component of $\eta$, and, for every
$e$ satisfying $\texttt{SideRead}$,
\[
 \begin{split}
  \kappa\in M.J&\ \Longleftrightarrow\
       \operatorname{loc}_s(\kappa)\in\operatorname{UniRT}(r),\\
  e(\chi(\kappa))&=
       a_s\!\left(\sum_{a=0}^{r-2} A_{\kappa,a,a+1}\right)
       \xi_s(\det D_\kappa).
 \end{split}                                                   \tag{AR.10}
\]
The formulas apply to $\kappa$ supported at $s.v$, with the character
formula restricted to the level. In $\operatorname{UniRT}(r)$ the
parahoric condition is supplemented by upper-unitriangular reduction
of the first $r$-by-$r$ block $A_\kappa$.

Define the residue additive character by
$z\mapsto\iota(\theta_a(\xi_{\mathrm{red},a}z))$.
Its value at $1$ is not $1$ when $r\ge2$. Insert the element with a
single $1$ on the first superdiagonal at this type place and identities
elsewhere into $\texttt{Bs.psiU}$. The formula for $\psi_U$ becomes
$\theta_a(\xi_{\mathrm{red},a})$, whereas
$X.\mathrm{psiU\_gen}$ gives the nontrivial generator value of its
nondegenerate additive character. Injectivity of the coefficient map
proves the assertion. AR6 extends the prescribed character.
For the formally possible $r<2$ case, choose any nontrivial residue
character; the superdiagonal sum is empty and no prescribed value
is used.

AR7 and AR1 give $\xi_s$ whose unit values are the fixed
$\tau_{j+1,a}$ residue values. Neither choice involves $e$.
AR2 makes all conductor and other type matrices the identity.
Unfolding $M.J$ therefore leaves exactly $\operatorname{UniRT}(r)$.
Its upper block reduces to an upper-unitriangular matrix, whose
determinant is $1$; the lower block has unit determinant by invertibility
of the full reduced matrix. The literal $\mathrm{gVal}$ expression
is the product of the additive superdiagonal character and the
$\tau$-character of that determinant. Reduction commutes with the
finite superdiagonal sum and with the determinant. The
$\texttt{SideRead}$ equality, $\texttt{Bs.psiU}$ and
$\texttt{TauTrue}$ now give (AR.10). This proves
$\texttt{levelUni\_export}$ with its original $\exists a_s,\xi_s\;\forall e$
quantifier order.

\paragraph{AR9: an identity frame identifies the two integral reduction operations.}
Let $P_V$ be the supplied $\texttt{IdFrame}$, with its exact rank
identification $\operatorname{rank}V=n$ and $D.\mathrm{host}.F=E^+$.
Let $\Phi=\mathrm{frameEquiv}(P_V)$. Then
$\kappa\in U_{\max}$ if and only if its framed matrix lies in the
stabilizer of the standard integral adelic lattice. If $\kappa$ is
integral, its framed residue entry is $\bar u$ exactly when the
corresponding selected local entry differs from $s.i(u)$ by an
element of valuation less than $1$.

The frame equation identifies each entry with the same finite adele
under the specified tensor interchange and rank reindexing. An
invertible matrix belongs to the integral general linear group
exactly when its entries and those of its inverse are integral:
the two matrices over the integral ring multiply to $1$, as is
checked after its injective inclusion in the adelic ring. This proves
the stabilizer assertion in both directions.
For the reduction assertion choose an integral $b$ representing the
global adelic residue. By construction the projected adelic entry
differs from $s.i(b)$ by an element of the maximal ideal. It differs
from $s.i(u)$ by such an element if and only if
$s.i(u-b)$ is in that ideal, by the ultrametric triangle inequality
in both directions. This last condition is exactly $u-b\in P$.
Thus the two residue operations coincide. Entrywise, this proves
the parabolic-block and upper-unitriangular equivalences used above
and below; there is no additional freedom to conjugate the frame.

\paragraph{AR10: a common unipotent datum gives the exact power comparison.}
Consider two geometric level data $L,L'$ in this same identity frame.
Assume the actual $\texttt{UniDat}$ data: their leading rank is $r$,
their blocks are $(r,t)$, their type-prime sets are the indexed $P_a$,
their conductor primes are split, lie in $\operatorname{range}P_q$
and avoid the type rational primes, and their residue-group
isomorphisms identify both level subgroups with the same upper
unitriangular group $U$. Finally assume their unipotent characters
are the same fixed function $f:U\to\mathbb C^\times$.
Then both levels contain every supported parahoric element at a
selected conductor place outside the type primes and every supported
$\operatorname{UniRT}(r)$ element at a type place. For an element
$\kappa$ in their intersection supported at $P_a$,
\[
 \chi_L(\kappa)^{g_a-1}
       =\chi_{L'}(\kappa)^{g_a-1}.                       \tag{AR.11}
\]

The containment statements follow from AR9. At other local places
the reductions are identity matrices. At the selected place the
block condition is exactly the required one; membership in $U$ is
detected after the specified residue-group isomorphism.
For a general element in the intersection, pull the two residue
tuples back to $U$. The entrywise equations for those isomorphisms
show that both preimages have the same matrix entries, hence are
the same $u\in U$. Both unipotent factors are therefore $f(u)$.
For an element supported at $P_a$, every remaining character factor
is $1$, except the diagonal block factors at $P_a$. Those factors
are characters of $\mathbb F_{g_a}^{\times}$ and their
$(g_a-1)$st powers are $1$. Raising the two products to that power
leaves the same $f(u)^{g_a-1}$, proving (AR.11).

\paragraph{AR11: the actual Joint is an intersection of two levels.}
Retain the precise input objects of $\texttt{joint\_of}$:
$Y:\texttt{R3.Read}$, $M_{\rm glob}:\texttt{MtFam}$,
$\texttt{MtAd}(Y,M_{\rm glob})$, the rational datum
$\texttt{RatAd}(Y,c_G,k_{K_s},I)$, its $\texttt{SideB}$, the
series $\texttt{SerAt}$ on that same datum and its same identity
frames, $\texttt{Keep}$, the actual level-character law,
$\texttt{Elem}$, and the selected-prime condition
$R.Q(t)\subset\operatorname{range}P_q$.
For $\mu$ satisfying $\texttt{Used0}(D,\mu)$ assume the exact equality
$\eta=R.C.\mu_0(j)\,R.\xi_H(Y.\mathrm{tupT}(\mu))$.
These data construct $\texttt{Joint}(\eta,M_{\rm glob}(\mu,j+1),X,b,h_A,j)$.

Let $L$ be the geometric level $\mathrm{datVP}$ with its literal
false-side level, and $L'$ the series level
$(\rho.B(i,j,b)).\mathrm{levS}(\rho.\mathrm{tw}(i,b,k,Q,\psi))$.
The residue isomorphisms come from $\texttt{Keep.grp}$, with
inverse $c_G$ on the geometric side. The rational law
$\texttt{RatAd.lvl}$ identifies its subgroup with $U$.
Its character law is
$\psi_U=(\psi_0)^{-1}$; the conjugated level in false-side
$\mathrm{datVP}$ contributes the other inverse. Thus its actual
unipotent factor is $f(u)=\iota(\psi_U(u))$.
On the series side the $\texttt{CellTie.U\_eq}$,
$\texttt{psiU\_eq}$, $\texttt{Bs.psiU}$ and $\rho.\mathrm{e\_LL0}$
fields give that same $f$. The selected-prime and block fields give
all remaining $\texttt{UniDat}$ clauses of AR10.

Take the frame automorphism to be the identity and set
\[
 J_1=\bigl(\mathrm{Lev.finE}\circ P_V.\psi\bigr)^{-1}
                   \bigl(\operatorname{lev}(L)\cap\operatorname{lev}(L')\bigr).
                                                               \tag{AR.12}
\]
Take $e=\rho.\mathrm{ek}(i,b,k)$, the polynomial target of the
series, and the injective semilinear map $\omega$ supplied by
$\texttt{SerAt.join}$. Its scalar restriction is the required side
embedding: equality on the coefficient ring extends to its fraction
field by writing each element as a fraction.
The exact join identity is
\[
 \mathrm{rdV}(f)(\Phi(g))=\omega((M.\mathrm{real}\,f)(g)). \tag{AR.13}
\]
AR10 proves both local containment fields for $J_1$. For
$\kappa\in J_1$ take the two scalars to be
$\chi_L(P_V.\psi(\kappa))$ and $\chi_{L'}(P_V.\psi(\kappa))$.
The $\texttt{MtAd.read}$ identity and injectivity of its coefficient
map transfer the first right-translation law from the actual
$Y.\mathrm{IL.jMt}$ form. For the second, (AR.13) transfers the
right-translation law of the series form. That law holds first on
the actual integral $\mathrm{MtLat}$ and then on its full fraction-field
span by linearity; it is not inferred merely from injectivity of
$\mathrm{rdV}$. AR10 proves the required equality of powers of the
two scalars. These are every field of $\texttt{Joint}$.

For the supplied stage, $\mu_{\mathrm{Tw}}(v_0)=\mu_{\mathrm{At}}(v_0)$
on the twisted support proves $\texttt{Used0}$. The literal
$\texttt{lev.xiH}$, $\texttt{tup}$ and $\texttt{tupTie}$ equations
give the displayed equality of $\eta$. Thus this also proves
$\texttt{joint}$ for that same stage. AR21 below constructs the
nontrivial reading used in (AR.13) for an installed series; the other
series and global realization fields are not created by (AR.12).

\paragraph{AR12: matrix intertwiners descend, and dual polynomiality is preserved.}
Let $e:K\hookrightarrow\mathbb C$, and let $A_\gamma,B_\gamma$ be
finite matrices over $K$, indexed by an arbitrary group. If a nonzero
complex matrix satisfies
\[
 X\,e(A_\gamma)=e(B_\gamma)\,X\qquad\hbox{for every }\gamma,
                                                               \tag{AR.14}
\]
then a nonzero matrix over $K$ satisfies the corresponding equations.
Choose a nonzero entry of $X$ and a $K$-linear functional
$\lambda:\mathbb C\to K$ that does not kill it; extend that entry
to a $K$-basis to obtain such a functional. Apply $\lambda$ entrywise.
It commutes with sums and multiplication by coefficients in $K$,
so every equation (AR.14) survives. The chosen entry proves that the
resulting matrix is nonzero. This works for infinitely many
$\gamma$ without an unjustified finite intersection argument.

There is also an explicit polynomiality statement needed in its
application. Suppose the dual of a finite-dimensional $K$-representation
of the rational unitary group has a basis in which its matrices are
polynomials in the entries of the embeddings of rational matrices.
Choose a basis of the original space and use its dual basis. If $A,B$
are the two constant change-of-basis matrices, then the original
matrix at $\gamma$ is the transpose of
$A\,\rho^\vee(\gamma^{-1})B$. On a unitary matrix $g$,
\[
 (g^{-1})_{ab}=\bar g_{ba}.
 \tag{AR.15}
\]
Replace each variable $(\sigma,a,b)$ in the polynomial matrix for
the dual by $(\sigma\circ c,b,a)$, where $c$ is the actual CM
involution, and perform these constant matrix multiplications and
the transpose. This gives polynomial matrices for the original
representation. Mapping their coefficients through $e$ proves
the literal $\texttt{IsAlgWeight}$ condition after complexification.
The source's $\texttt{HWAtF}$ includes the required polynomial
matrix presentation; existence of a weight vector alone would not
suffice for this argument.

\paragraph{AR13: the polynomial coefficient representation descends over the actual field.}
We state precisely the representation-theoretic inputs retained here.
First, an irreducible algebraic representation of the rational definite
unitary group, for a complete CM type, is isomorphic over $\mathbb C$
to the canonical tensor coefficient module $V_{\mathbb C}(\lambda)$
for dominant integral weights $\lambda$. Second, the dual of that
canonical module has exactly the highest weights specified in
$\texttt{contragredient\_polynomialWeight\_unique\_highestWeight}$,
with the source's conjugate-CM-type convention. These are the named
classification and canonical highest-weight inputs, not conclusions
of finite-dimensionality.

Under these two inputs, let $\rho$ be an irreducible finite-dimensional
representation over $K$ and suppose its dual has the literal
$\texttt{HWAtF}$ weights $\lambda_0(\tau)$ for each embedding in a
complete CM type. Then there are dominant weights $\lambda$ and
a $K$-linear equivariant isomorphism
\[
 T:W\xrightarrow{\ \sim\ }V_K(\lambda),\qquad
             \lambda(\tau)=\lambda_0(\tau).             \tag{AR.16}
\]
In particular the field of definition is $K$, not a newly chosen
splitting field.

Choose a $K$-basis of $W$ and complexify its matrices using $e$.
AR12 proves algebraicity. A nonzero complex representation contains
an irreducible subrepresentation: choose a nonzero invariant subspace
of minimal dimension. Algebraicity passes to this subspace: express
its inclusion by a constant matrix, choose a constant left inverse,
and multiply the ambient polynomial matrix by these two matrices.
By the stated classification this subspace is
$V_{\mathbb C}(\lambda)$ for some $\lambda$. Its inclusion in
$W_{\mathbb C}$ is a nonzero complex intertwiner. The canonical
coefficient modules have bases $b_K,b_{\mathbb C}$ on the same finite
index set, with $b_{\mathbb C}=e(b_K)$, and their action matrices
commute with coefficient extension. AR12 descends the inclusion to
a nonzero $K$-intertwiner
$\alpha:V_K(\lambda)\to W$.
Its image is a nonzero invariant subspace, so irreducibility of
$W$ makes $\alpha$ surjective. On the other hand the original
complex inclusion gives
$\dim V_K(\lambda)=\dim V_{\mathbb C}(\lambda)\le\dim_K W$.
Surjectivity gives the reverse inequality. Thus $\alpha$ is an
isomorphism; there is no need to assume that complexification of an
arbitrary irreducible representation stays irreducible.
Transfer the prescribed dual highest-weight equations across
$\alpha$ and then across the common coefficient bases. The second
named input gives $\lambda_0(\tau)=\lambda(\tau)$.
This proves $\texttt{irreducible\_weight\_descends\_to\_polynomial\_module}$
from the two precise complex inputs.

\paragraph{AR14: the actual infinitesimal parameters give CellWt.}
For antitone integer sequences $a,b$ of length $n$, equality of
their source infinitesimal-character multisets implies $a=b$.
Indeed the entries are
$a_i-i+(n-1)/2$. The sequence $a_i-i$ is strictly decreasing, and
so is $b_i-i$. Removing the common shift and comparing successively
the largest entries of the two finite multisets proves equality
at each index. This proves $\texttt{eq\_of\_infOf\_eq}$, including
the case $n=0$.

Apply the complex classification input of AR13 to the actual
$m.\mathrm{wt}$ in $\texttt{MtReal}$, obtaining weights
$\lambda_C$. Its supplied $\texttt{Mt}$ weight field, transferred
through the canonical duality with the conjugate CM type, says
\[
 \operatorname{infOf}(\lambda_C(\tau))
           =\operatorname{paramV}(R,hc_\tau,j,1/2).
                                                               \tag{AR.17}
\]
For the coefficient representation $X.\rho_W(i,j+1,b)$, its
$\texttt{rhoW\_hw}$ field supplies $\lambda_0$.
The exact $\texttt{CharSide}$ assertion
$\operatorname{infG}(\eta_T)=\{1/2\}$ and
$\texttt{rhoW\_wt}$ give the same right side of (AR.17).
AR13 supplies dominant $\lambda_K=\lambda_0$ and an isomorphism
over $K_0=\operatorname{Frac}(O_i)$. The decreasing-sequence argument
therefore gives $\lambda_K=\lambda_C$.
The equality
\[
 (e\circ\mathrm{toKk})\circ\mathrm{kEmb}_\tau
                      =\mathrm{embT}_\tau
                                                               \tag{AR.18}
\]
is exactly $\texttt{SideRead}$ applied to the specified embedded
host element. The actual CM type is injective, disjoint from its
conjugate and covers the complex embeddings, as required in AR13.
The host conjugation is the CM involution because their images
under the specified complex embedding agree.

Pulling back the common canonical polynomial bases through the two
isomorphisms now gives $\texttt{WtMatch}$ over $e\circ\mathrm{toKk}$.
For clarity, if $b_K$ is a canonical basis, then
\[
 \operatorname{coord}_{e(b_K)}(e(y))
                         =e(\operatorname{coord}_{b_K}(y)),
                                                               \tag{AR.19}
\]
by expanding $y$ in that basis. The polynomial action commutes with
coefficient extension, so (AR.19) proves entrywise equality of
every representation matrix.

Finally the actual cell field $\texttt{bc}$ supplies the base-change
matrices from $K_0$ to $K_{i,k}$. Its basis need not be the one just
chosen. If $P,Q$ are inverse change-of-basis matrices over $K_0$,
their images remain inverse over $K_{i,k}$; change the cell basis
by $P$ and use $QAP$ for the action matrices. Applying $e$ gives
the same conjugated complex matrices. Thus
\[
 \texttt{WtMatch}(e,M.\rho,m.\mathrm{wt}.\rho)
                      =\texttt{CellWt}(m,X,b,h_A,j,e)            \tag{AR.20}
\]
holds for every allowed $e$. This proves $\texttt{cellWt\_all}$
under the actual weight fields and the two named classification
inputs; it does not construct those weight fields or $\texttt{MtReal}$.

\paragraph{AR15: the coefficient module is the span of its extremal orbit over a field.}
Here $V_K(\lambda)$ is the tensor of the source's \emph{dual Weyl}
coefficient modules, with determinant shifts. We retain the standard
characteristic-zero equality of the divided-power Weyl module and
the dual Weyl module over $K$, the literal
$\texttt{dualWeyl\_eq\_weyl}$. No integral equality is inferred.
Under that input,
\[
 V_K(\lambda)
   =\operatorname{span}_K
       \{\rho_{\mathrm{poly}}(g)\Delta_\lambda:
                         g\in\prod_\tau\mathrm{GL}_n(K)\}.
                                                               \tag{AR.21}
\]

The nontrivial direction is elementary interpolation. A subspace
stable under row substitutions and containing a polynomial $F$
contains every coefficient of its substitution
$F(X(I+zE_{ba}))$. For each nonnegative integer $z$, the matrix is
invertible: if $a\ne b$ it is a transvection, and if $a=b$ its
nontrivial diagonal entry is $1+z\ne0$ in characteristic zero.
Taking more distinct such $z$ than the degree and inverting the
Vandermonde matrix expresses every coefficient as a linear
combination of these substitutions. Those coefficients are exactly
the row divided powers. By iteration the orbit span contains every
divided-power word applied to the highest polynomial.
The retained characteristic-zero equality makes these span the
dual Weyl module. Tensoring the argument over the embeddings and
absorbing the nonzero determinant characters proves (AR.21).
The reverse inclusion follows directly from stability of that
coefficient module under the polynomial action.
This assertion concerns fields of characteristic zero. It does
not identify the integral lattices $\Delta$ and $\nabla$ in CS.

\paragraph{AR16: complementary minors give the harmonic extremal value and its full image.}
Let $w$ be the actual good weight, with positive and negative
lengths $p_\tau,q_\tau$ satisfying $p_\tau+q_\tau\le m_F$ and
$p_\tau+q_\tau\le n$. Here $m_F$ is the Fock row rank; it is not the
global chain length. Write $a_\tau^+,a_\tau^-$ for its decreasing
nonnegative row lengths and put
$d_{\tau,i}^{\pm}=a_{\tau,i}^{\pm}-a_{\tau,i+1}^{\pm}$, extending
the last row length by zero. Let $\Delta_F=w.\delta(m_F)$, and let
$\mathrm{ev}_y$ evaluate at the source's $\mathrm{fockPt}(w.p,y)$.
For the leading principal minor $L_k(y)$,
\[
 \begin{split}
 \mathrm{ev}_y(\mathrm{minorPlus}_{\tau,k})&=L_k(y_\tau),\\
 \mathrm{ev}_y(\mathrm{minorMinus}_{\tau,k})
    &=(-1)^{k(k-1)/2}\,
                L_{n-k}(y_\tau)\det(y_\tau)^{-1}.
 \end{split}                                                   \tag{AR.22}
\]

The first equality is substitution into a determinant. For the
second, substitution produces the bottom principal minor of the
inverse matrix, with one row reversal. The complementary-minor
identity changes that minor to $L_{n-k}/\det$, while the reversal
has sign $(-1)^{k(k-1)/2}$. The identity holds on all invertible
matrices: multiply by the determinant and use the adjugate
cofactor identity, so no invertibility of an intermediate leading
minor is required.

Multiplying (AR.22), with exponents $d^\pm$, gives
\[
 \mathrm{harmFn}_{B}(\Delta_F)(y)
       =s_w\,\mathrm{coefFn}_{\lambda}(\Delta_\lambda)(B^{-1}y),
 \quad
 s_w=\prod_{\tau}\prod_{i=0}^{q_\tau-1}
             (-1)^{i(i+1)d^-_{\tau,i}/2}\in\{1,-1\}.
                                                               \tag{AR.23}
\]
The exponents agree because consecutive drops telescope:
$\sum_i d^-_{\tau,i}=a^-_{\tau,0}$, and the exponent of each
proper leading minor is
$\lambda_{\tau,k}-\lambda_{\tau,k+1}$; the remaining determinant
exponent is $\lambda_{\tau,n-1}$. This also covers an empty positive
or negative block. It verifies the exact sign, not just a nonzero
proportionality.

The space $\mathrm{harm}(m_F,w,B)$ is, by definition, the span of
the Fock-group orbit of $\Delta_F$. Direct substitution gives
$\mathrm{harmFn}(g\cdot P)(y)=\mathrm{harmFn}(P)(g^{-1}y)$.
The coefficient action has the same law after the inverse-transpose
change $g\mapsto g^{-T}$. As $g$ varies this is a bijection, and
$s_w$ is nonzero. AR15 and (AR.23) therefore show
\[
 \mathrm{harmFn}_{B}(\mathrm{harm}(m_F,w,B))
              =\mathrm{coefFn}_{\lambda}(V_{\mathbb C}(\lambda)).
                                                               \tag{AR.24}
\]
This proves the harmonic image identity, conditional only on the
explicit characteristic-zero input in AR15.

\paragraph{AR17: evaluation is injective on the harmonic orbit space.}
The harmonic evaluation map in (AR.24) is injective. This assertion
has a direct positive-inner-product proof and does not require
the classification input of AR13.
On the polynomial algebra use the Fock pairing
\[
 \langle X^\alpha,X^\beta\rangle_F
                       =\delta_{\alpha\beta}\,\alpha!.
                                                               \tag{AR.25}
\]
It is positive definite. For the mixed Fock substitution $T(g)$,
the plus variable with column index $a$ maps to
$\sum_b(g^{-1})_{ab}X_b$, whereas the minus variable maps to
$\sum_b g_{ba}X_b$, with the row and embedding indices fixed.
Applying the multinomial
formula to monomials proves
$T(g)^*=T(g^*)$, where $g^*$ is conjugate transpose in each factor.
Both sides stay in the finite homogeneous polynomial space in use.

Let $e_0=\mathrm{ev}_1(\Delta_F)$ and
$A=\langle\Delta_F,\Delta_F\rangle_F$.
Formula (AR.22) at $1$ gives $e_0=s_w\ne0$.
For every $g$,
\[
 e_0\langle\Delta_F,T(g)\Delta_F\rangle_F
                         =A\,\mathrm{ev}_1(T(g)\Delta_F).       \tag{AR.26}
\]
To verify it first take $g=ul$, with $u$ upper unitriangular
and $l$ lower triangular. Each leading flag minor shows that
$T(l)\Delta_F=c(l)\Delta_F$, with $c(l)=1$ when $l$ is lower
unitriangular. In particular $T(u^*)\Delta_F=\Delta_F$, so
adjointness gives
$\langle\Delta_F,T(u)\Delta_F\rangle_F=A$.
Evaluating the minors at the upper unitriangular $u^{-1}$ gives
$\mathrm{ev}_1(T(u)\Delta_F)=e_0$.
Multiplication by $c(l)$ proves (AR.26) on this big cell.
Every invertible matrix can be approached by matrices in that cell:
for $(1-z)g+zI$, each required trailing principal minor is a
polynomial in $z$ that is nonzero at $z=1$. Avoid their finitely many
zeros and let $z\to0$. Substitution and inverse matrices are
continuous at $g$. Thus (AR.26) holds everywhere.

By linearity, (AR.26) holds for every polynomial in the orbit span.
If $P$ is in that span and its harmonic function is zero, then
$\mathrm{ev}_1(T(g)P)=0$ for every $g$; the determinant prefactor
in $\mathrm{harmFn}$ is never zero. Equation (AR.26) gives
$\langle\Delta_F,T(g)P\rangle_F=0$.
Adjointness then says
$\langle T(g)\Delta_F,P\rangle_F=0$ for every $g$.
Because $P$ itself is in their span, $\langle P,P\rangle_F=0$,
and positivity proves $P=0$.

The coefficient-function map on $V_{\mathbb C}(\lambda)$ is also
injective: remove its nonzero determinant factor and obtain a
polynomial vanishing on the product of general linear groups.
That set is a nonempty Zariski-open dense subset of the matrix
space, so the polynomial is zero; alternatively apply the
one-variable finite-zero argument successively to its entries.
Together with AR16 this gives a unique linear isomorphism
\[
 \Psi:V_{\mathbb C}(\lambda)\xrightarrow{\ \sim\ }
                      \mathrm{harm}(m_F,w,B),
 \qquad \mathrm{harmFn}\circ\Psi=\mathrm{coefFn}.        \tag{AR.27}
\]
Its equivariance follows by applying the injective harmonic
evaluation map to the two candidate translates.

\paragraph{AR18: the coefficient embedding has all five WtMap properties.}
Retain a good geometric base $B$, its identity frame, the actual
rank equality, a coefficient ring $R$ with fraction field $K$,
$e:K\hookrightarrow\mathbb C$, $f:K_0\hookrightarrow K$, and
$s:L_0\hookrightarrow K_0$ with $efs=\iota$.
Assume the literal $\texttt{IsCoefRing}$, the common weight
$\lambda=B.w.\lambda$, an equivariant
$T_0:W_0\simeq V_{K_0}(\lambda)$, and the actual base-change
matrices for $W$ over $K$. Then there is an injective
$e$-semilinear map $\omega:W\to\mathrm{PolF}$ satisfying every
field of $\texttt{WtMap}$.

Choose common base-change bases $b_0,b_1$, set $c_0=T_0b_0$, and define
\[
 \Omega(w)=\sum_a e((b_1^*w)_a)\,
                      (ef)(c_{0,a}),\qquad \omega=\Psi\Omega.    \tag{AR.28}
\]
Images of the $c_{0,a}$ form a complex basis of the canonical module:
change from its canonical coefficient basis and apply the field
embedding to the inverse change-of-basis matrices. Thus $\Omega$,
and hence $\omega$, is injective. Its values are harmonic by (AR.27).
The same matrix equation used in AR14 shows that it intertwines
the rational unitary actions. The identity frame identifies the
rational matrices on both sides, so this is the literal
geometric action in $\texttt{WtMap}$.

Two lattice assertions remain, and both are essential. Let
$O=e(R)\subset\mathbb C$ and let $\Lambda_h$ be the geometric
coefficient lattice at the actual point $h$. By definition it
consists of harmonic polynomials whose harmonic functions belong
to the $O$-span of the coefficient functions of
$\nabla_{\!L_0}(\mathrm{latAt}(h))$. Each such generator lifts through
$s$ to $V_{K_0}(\lambda)$, and therefore through (AR.28) to $W$.
Its $O$-linear combinations also lift, using $R$-scalars.
Injectivity of coefficient evaluation shows that the resulting
image is exactly the desired harmonic polynomial. Thus
$\Lambda_h\subset\operatorname{range}\omega$.

Conversely, for every $w$ and $h$ there is a nonzero $c\in R$ with
$\omega(cw)\in\Lambda_h$. The lattice
$\nabla_{\!L_0}(\mathrm{latAt}(h))$ spans the coefficient module.
One way to check this is to sandwich the underlying full lattice
between two nonzero integral multiples of a standard lattice,
apply the finite tensor polynomial construction, and clear the
determinant shifts. Its coefficient-field span is unchanged.
Express the finitely many $c_{0,a}$ in the image of that spanning
lattice. Their finitely many coefficients, multiplied by the
coordinates of $w$, lie in $K=\operatorname{Frac}(R)$.
Clear their denominators and take the product of the nonzero
denominators. This gives the required $c$.
These are precisely the range and denominator fields of
$\texttt{WtMap}$, in addition to injectivity, harmonic membership
and equivariance. No surjectivity onto all complex harmonic
polynomials over the smaller field is asserted.

\paragraph{AR19: the global level equality includes every character factor.}
Assume the actual setup $\rho$ is good, pinned to the same
$R,c_G$, has the prescribed level of $D.\mathrm{nexp}$, and has
choice set $\operatorname{range}P_q$ at split primes.
Retain $\texttt{Keep}$, $\texttt{TauTrue}$, $\texttt{ValLL0}$,
the actual base-character law and the stated split auxiliary
places. Then, for every actual $\texttt{TwOK}$ twist,
\[
 g\in M.J\ \Longleftrightarrow\ \Phi(g)\in B.\mathrm{levJ}(t),
 \qquad
 e(\chi(g))=B.\mathrm{levLam}(t)(\Phi(g)).                \tag{AR.29}
\]
This is $\texttt{LevMatch}$, not merely equality on elements
supported at a single place.

AR2 identifies the two selected conductor sets exactly, and the
pinning gives the same type primes, residue isomorphisms and
identity frame. AR9 converts the maximal-compact and every
block-reduction condition. The entrywise residue-group equation
identifies membership of the type tuple in $U$, proving the first
part of (AR.29) for an arbitrary $g$.

For its character, use the actual product formulas on both sides.
The unipotent character is a homomorphism on the product of the
type groups. Decomposing a tuple into its single-place tuples
expresses its value as the product of the corresponding values;
the upper-unitriangular hypothesis ensures each tuple is in its
domain. At each type place, AR7 reads the diagonal residue
character, and the common $\psi_U$ identity reads the unipotent
factor. At $Q_P$, AR3 reads the residue character of $\psi$; at
$D_P$, AR4 reads the residue character of $\delta_i^{\epsilon_b}$.
For any integral invertible matrix satisfying the block condition,
reduction of the lower-right determinant equals the determinant
of the reduced lower block. Both are units, as in AR5.
Hence the same integral representative can be used in every
corresponding factor. Multiplication proves (AR.29).
On the false side the auxiliary Hecke character is first the
character of $\delta_i^{-1}$ and the residue exponent is $-1$;
these operations must be kept together. The calculation gives
the actual cell factor in (AR.5), rather than an extra inverse.

\paragraph{AR20: the actual form reading is pointwise and its integral preimage spans.}
Here is the complete linear-algebra construction used by
$\texttt{read\_cell}$. Suppose $\Phi:G\to G'$ is a group isomorphism,
the rational subgroups correspond in both directions, and
the weight and level comparisons are those of AR18 and AR19.
Let
$\mathrm{real}:V\simeq\mathcal A_K(G,\Gamma,J,\rho,\chi)$
be the supplied actual realization. Assume finitely many
representatives $y_a$ for the geometric double cosets and the
actual covariance of the coefficient lattices. Define
\[
 r(f)(h)=\omega((\mathrm{real}\,f)(\Phi^{-1}h)).          \tag{AR.30}
\]
Then $r$ is injective and semilinear, every integral geometric
form is in its range, and
\[
 \operatorname{span}_K\{f:r(f)\in B.\mathrm{mtLat}\}=V.  \tag{AR.31}
\]

The two transformation laws for $r(f)$ follow by substitution in
(AR.30), rational-subgroup correspondence, and (AR.29).
Injectivity follows pointwise from injectivity of $\omega$ and
$\mathrm{real}$. For an integral geometric form $F$, AR18 gives
a unique preimage $u(g)\in W$ of each $F(\Phi(g))$.
Apply $\omega$ to the required left and right covariance
equations for $u$. They become the known equations for $F$, so
injectivity proves them. Thus $u$ lies in the source form space
and $\mathrm{real}^{-1}(u)$ is the desired preimage.

For an arbitrary $f\in V$, choose a nonzero denominator $c_a\in R$
at each representative $y_a$ using the fifth $\texttt{WtMap}$ field.
Put $c=\prod_a c_a\ne0$. At $y_a$, the value of $r(cf)$ is
$(c/c_a)$ times the integral value provided by $c_a$. Here
$c/c_a=\prod_{b\ne a}c_b$ belongs to $R$, so this value is in the
coefficient lattice. Every point is $\gamma y_a\kappa$.
The geometric action transports that lattice to the required
one, and the right level scalar lies in $e(R)$. To see the last
assertion, the level scalar is a product of characters of finite
residue groups, hence a root of unity. Its preimage under $e$ is
integral over $R$, being a root of $X^N-1$. The actual coefficient
ring is integrally closed, so that preimage belongs to $R$.
Thus $r(cf)$ is integral everywhere. Finally $f=c^{-1}(cf)$ proves
(AR.31). This proof does not require the representatives to have
trivial stabilizer, since the pre-existing form already has its
stabilizer covariance.

\paragraph{AR21: the installed reading supplies the nontrivial series join.}
Under the full input list of
$\texttt{ExistsReadingRealSide.exists\_reading}$, AR13 and AR14
identify the coefficient weight with $B.w.\lambda$:
the setup's good-base weight field and the pinning give the
literal lift weight, and equality of its infinitesimal parameter
uses the same decreasing-sequence argument.
AR18 constructs $\omega$, AR19 gives $\texttt{LevMatch}$, and AR20
constructs $r=\mathrm{rdV}$ with all three required properties:
injectivity, range containing the integral forms, and full span
of their preimage. It also gives, by its definition,
\[
 r(f)(\Phi(g))=\omega((M.\mathrm{real}\,f)(g)),
                                                               \tag{AR.32}
\]
the exact join equation.

The choices are made once for each $(i,j,b,k,Q,\psi)$. If
$\texttt{TwOK}$ holds, make the construction above; if it fails,
choose the zero map, because the conclusions are all conditional
on $\texttt{TwOK}$. This preserves the source's quantifiers and
does not require all possible twists to be admissible. Proof
arguments witnessing admissibility do not alter the cell data.
When $\texttt{exists\_tie\_withSeries}$ installs this reading, its
explicit equality $\mathrm{CellTie.rdV}=r$ converts (AR.32) into
$\texttt{SerAt.join}$. Changing between two identity frames does
not change the finite adelic map: their entrywise matrix-preservation
equations are the same, and the ambient embeddings are injective.
This is exactly the final frame replacement in
$\texttt{exists\_preTie}$.

Thus the join is supplied by an actual coefficient map and
denominator argument for the installed series. We have not proved
the whole $\texttt{exists\_preTie}$ theorem from arithmetic data:
its good pinned setup, scalar normalization, faithful reading,
theta/class-sum installation, Hecke tie and universal-character
fields remain their separately stated inputs. Nor may an arbitrary
previous $\mathrm{rdV}$ be silently replaced by (AR.30).

\paragraph{AR22: the precise rank and choice boundary after these constructions.}
For the same supplied stage, AR1, AR5 and AR8 provide the missing
bad-place and local-shape inputs in RC19. AR11 supplies its actual
$\texttt{Joint}$, with AR21 available for the installed reading.
The existing product-level and nonzero-form arguments in RC14--RC19
therefore give the literal $\texttt{LevEq}$. Independently AR14
gives $\texttt{CellWt}$. RC1--RC2 now give the required dimension
equality and normalized $\texttt{RankTie}$ for each twist under
these hypotheses. This completes the local arithmetic-comparison
part of $\texttt{StatementCOfCasselmanWallach}$ with its specified
global realizations and coefficient structures.

This conclusion does not manufacture $\texttt{LocsCarry}$.
For every actual local index its full
$\texttt{IsGalOf}(D,a,\mathrm{univ},\psi_{L,a},\chi_{w(a)},0)$,
its selected $\mathrm{twOf}$ and its side identities are still needed,
as in RC10. With these, the nonvacuous all-place twist readings,
and the remaining exact choice hypotheses, NC and RC11 construct
the normalized $\texttt{NChoice}$. The separate
$\texttt{MuIsW}$ join still needs the same $\xi,\psi,Q$ and the
actual $\texttt{ColW}$, $\texttt{Pre2}$, $\texttt{RowPk}$ and
$\texttt{St}$ identities of RC13.
The bare LS stage has empty twist and local index sets and cannot
supply $\texttt{Bs.cover}$ in AR7 or the all-place twist readings
required by this construction.

The remaining substantive suppliers are therefore explicit:
global reciprocity and the normalized algebraic adic-character
constructor; the two complex highest-weight inputs and the stated
characteristic-zero coefficient-module identity; compatible
nonempty twist/local data with their normalized character and
weight fields; and the common global realization/pinning/series
installation. The complex HFrame character constructor addresses
its own part of these data. No analytic host/window existence,
local-to-global realization, or RH/GRH conclusion follows merely
from the conditional rank equality proved here.
